\documentclass[10pt,a4paper]{amsart}
\usepackage[margin=19mm]{geometry}
\usepackage{amsmath,amssymb,mathtools}
\usepackage[scr=rsfs,cal=euler]{mathalfa}
\usepackage{xfrac}
\usepackage[unicode,hidelinks,bookmarks=false]{hyperref}
\newcommand{\ie}{i.e.\ }
\newcommand{\cf}{cf.\ }
\newcommand{\resp}{respectively\ }
\newcommand{\Subsection}{\subsection}
\providecommand{\nobreakdash}{\nobreak\hbox{-}}
\setlength{\emergencystretch}{2em}
\multlinegap0pt
\usepackage{amssymb}
\usepackage{enumerate}
\usepackage[normalem]{ulem}
\newtheorem{lemma}{Lemma}
\newcommand\C{{\mathbb{C}}}
\newcommand{\ep}{\varepsilon}
\title{Convergence of higher derivatives of random polynomials with independent roots}
\author{J\"urgen Angst, Oanh Nguyen, Guillaume Poly}
\date{Source: arXiv:2601.01212v1 (CC BY 4.0)}
\begin{document}
\maketitle

\noindent\emph{Source and licence.} Convergence of higher derivatives of random polynomials with independent roots. Version arXiv:2601.01212v1 (CC BY 4.0), distributed by arXiv under the Creative Commons Attribution 4.0 International licence, http://creativecommons.org/licenses/by/4.0/, as recorded in the arXiv licence metadata for this exact version and shown on the abstract page https://arxiv.org/abs/2601.01212v1. Full licence notice: \texttt{licenses/arxiv-2601.01212-CC-BY-4.0.txt}.\par\smallskip

\noindent\textit{Editorial scope.} Counted unit: the article's lemma lm:splitting (source0.tex lines 528--532). The article's complete original proof of it is delivered with it (source0.tex lines 544--636, one \texttt{proof} environment of 93 lines, which the article places away from the statement). No mathematics is written, repaired or re-derived: the statement and the proof are the article's own lines, in the article's own notation, and no retained line is substituted or re-flowed.  Retained uncounted as the source-local context the proof needs: lines 276--278.  Nothing was omitted from any retained span, so the package carries no omission record.  No retained line contains a citation, so the wrapper carries no bibliography and no bibliography entry.  No retained line contains a figure environment, an image-inclusion command or drawing code, so no image is delivered and the package declares no asset dependency.  Editorial formatting only, in addition: an AMS \texttt{amsart} preamble that restates the article's own declarations and macros used by the retained lines, namely the environments \texttt{lemma} on separate counters, together with the standard packages those lines need.  The retained environments print in this document as Lemma 1 in order of appearance, whereas the article prints the same items with the numbering of its own full document.  The complete native source of the article as distributed in the arXiv e-print for this exact version is bundled unchanged as \texttt{sources/192063/source0.tex}.
 \begin{equation}
 	\mu(A)\ge c_0 \text{Leb}(A).\label{cond:doeblin}
 \end{equation}

 \begin{lemma}\label{lm:splitting}
 	Assume that the measure $\mu$ satisfies the Doeblin's condition \ref{cond:doeblin}, for all $z\neq z_0$, there exist parameters $c_a\in (0, 1), r_a>0, w_a\in \C$ such that  
 	$$Y_i = \frac{1}{z-\xi_i} \overset{d}{=} \ep_i Z_i + (1-\ep_i)W_i$$
 	where $\ep_i, Z_i, W_i$ are independent random variables with $\ep_i\sim \text{Ber}(c_a)$, $Z_i\sim \text{Unif}(B(w_a, r_a))$.
 \end{lemma}

 \begin{proof} [Proof of Lemma \ref{lm:splitting}]
 	Let
 	\[
 	Y := \frac{1}{a - \xi}, \qquad \xi \sim \mu.
 	\]
 	We first show that $Y$ also satisfies the Doeblin's condition, possibly with different parameters. Indeed, set $\delta := |a - z_0|>0$ and choose
 	\[
 	r_a := \min\{r_0, \delta/2\}.
 	\]
 	Then for every $w \in B(z_0, r_a)$ we have $|a - w|\ge \delta/2.$
 	Define
 	\[
 	T_a : B(z_0, r_a) \to \C, \qquad T_a(w) := \frac{1}{a - w}.
 	\]
 	The map $T_a$ is holomorphic on $B(z_0, r_a)$, with derivative
 	\[
 	T_a'(w) = \frac{1}{(a-w)^2}.
 	\]
 	Since on $B(z_0,r_a)$ we have
 	\[
 	\delta/2 \le |a-w| \le |a - z_0| + r_a \le \delta + \delta/2 = 3\delta/2,
 	\]
 	we obtain that $ |T_a'(w)|$ is uniformly bounded above and below. Therefore, $T_a$ is a $C^1$-diffeomorphism from $B(z_0, r_a)$ onto the open set
 	\[
 	D_a := T_a\big(B(z_0, r_a)\big).
 	\]
 	Moreover, the Jacobian of the inverse $T_a^{-1}$ is
 	\[
 	J_{T_a^{-1}}(y) = \frac{1}{J_{T_a}(T_a^{-1}(y))},
 	\]
 	so it is also bounded above and below on $D_a$.
 	
 	\medskip
 	
 	Let $\nu_a$ be the distribution of $Y$, then for any Borel set $E \subset D_a$, we have $T_a^{-1}(E) \subset B(z_0, r_a) \subset B(z_0, r_0)$, so from \eqref{cond:doeblin},
 	\begin{equation}
 		\nu_a(E) = \mu\big(T_a^{-1}(E)\big) \ge c_0 \, \mathrm{Leb}\big(T_a^{-1}(E)\big).
 		\label{eq:nu-lower}	 		
 	\end{equation}
 	Since $T_a$ is a diffeomorphism on this set and $J_{T_a^{-1}}$ is bounded below there exists $m_a>0$ (depending only on $a, z_0, r_0, c_0$) such that
 	\begin{equation}
 		\mathrm{Leb}\big(T_a^{-1}(E)\big)
 		= \int_{E} J_{T_a^{-1}}(y)\, dy
 		\ge m_a \int_{E} dy
 		= m_a \, \mathrm{Leb}(E)
 		\qquad \text{for all Borel } E \subset D_a.
 		\label{eq:jac-lower}
 	\end{equation}
 	Combining \eqref{eq:nu-lower} and \eqref{eq:jac-lower} we get
 	\begin{equation}
 		\nu_a(E) \ge c_0 m_a \, \mathrm{Leb}(E)
 		\qquad \text{for all Borel } E \subset D_a.
 		\label{eq:doeblin-Y-a}
 	\end{equation}
 	
 	\medskip
 	
 	Since $D_a$ is open, there exists a ball $B(w_a, r_a) \subset D_a$.  
 	Then from \eqref{eq:doeblin-Y-a},
 	\[
 	\nu_a(E) \ge c_0 m_a \, \mathrm{Leb}(E)
 	\qquad \text{for all Borel } E \subset B(w_a, r_a).
 	\]
 	Thus, $Y$ satisfies the Doeblin's condition. We then set
 	\[
 	c_a := c_0 m_a \mathrm{Leb}\big(B(w_a, r_a)\big) \le \nu_a(B(w_a, r_a)).
 	\]
 	Thus, $c_a\in (0, 1]$.
 	If $c_a=1$ the decomposition is trivial; otherwise we define the uniform measure
 	\[
 	\lambda_a(E) := \frac{\mathrm{Leb}(E \cap B(w_a, r_a))}{\mathrm{Leb}(B(w_a, r_a))}.
 	\]
 	Then for every $E \subset B(w_a, r_a)$,
 	\begin{equation}
 		\nu_a(E) \ge c_a \, \lambda_a(E).
 	\end{equation}
 	
 	\medskip
 	
 	Define the probability measure
 	\[
 	\eta_a := \frac{\nu_a - c_a \lambda_a}{1 - c_a}.
 	\]
 	Then we have the decomposition
 	\[
 	\nu_a = c_a \lambda_a + (1-c_a)\eta_a.
 	\]
 	Equivalently, we can realize $Y$ as
 	\[
 	Y \overset{d}{=} \ep Z + (1-\ep) W,
 	\]
 	where $\ep  \sim \mathrm{Ber}(c_a)$, $Z \sim \lambda_a$ , $W \sim \eta_a$, and $\delta_a, Z, W$ are independent.
 \end{proof}

\end{document}
